\documentclass[12pt]{article}
\usepackage[margin=1in]{geometry}
\usepackage{amsmath,amssymb,amsthm,algorithm,algorithmic}
\usepackage{enumitem}

\newtheorem{theorem}{Theorem}
\newtheorem{lemma}[theorem]{Lemma}
\newtheorem{definition}[theorem]{Definition}

\title{Problem 96: Su Doku}
\author{Project Euler Solution}
\date{}

\begin{document}
\maketitle

\section{Problem Statement}
Solve fifty Sudoku puzzles and compute the sum of the three-digit numbers formed by the first three cells of each solution's top row.

\section{Mathematical Foundation}

\begin{definition}[Sudoku as a CSP]
A Sudoku instance is a constraint satisfaction problem on $G = \{0,\ldots,8\}^2$ with variables $v_{r,c}$, domains $D_{r,c} \subseteq \{1,\ldots,9\}$, and 27 all-different constraints (one per row, column, and $3 \times 3$ box).
\end{definition}

\begin{definition}[Peers]
The \emph{peers} of $(r,c)$ are the $20$ cells sharing a row, column, or box with it.
\end{definition}

\begin{theorem}[Soundness of propagation]
\label{thm:prop}
If $v_{r,c} = d$, removing $d$ from every peer's domain preserves all valid completions.
\end{theorem}
\begin{proof}
In any valid completion, the all-different constraint on each unit containing $(r,c)$ forces $v_{r',c'} \ne d$ for every peer $(r',c')$. Hence $d \notin D_{r',c'}$ in any solution, and removal is safe.
\end{proof}

\begin{lemma}[Naked single]
If $|D_{r,c}| = 1$ after propagation, the sole remaining value is forced.
\end{lemma}
\begin{proof}
$v_{r,c}$ must take a value in $D_{r,c}$; a singleton domain admits exactly one choice.
\end{proof}

\begin{lemma}[Hidden single]
If value $d$ appears in exactly one cell's domain within a unit, that cell must hold~$d$.
\end{lemma}
\begin{proof}
The all-different constraint requires $d$ to appear in the unit. With only one possible location, the assignment is forced.
\end{proof}

\begin{theorem}[Completeness of backtracking]
\label{thm:complete}
Backtracking with propagation (Theorem~\ref{thm:prop}) is complete: it finds a solution if one exists.
\end{theorem}
\begin{proof}
The search tree has finite depth ($\le 81$) and bounded branching factor ($\le 9$). Propagation removes only values inconsistent with all valid completions, so no solution path is pruned. Exhaustive exploration of all surviving branches guarantees that a solution is found if it exists.
\end{proof}

\begin{theorem}[MRV heuristic]
Selecting the unassigned cell with minimum $|D_{r,c}|$ minimizes the branching factor at each node.
\end{theorem}
\begin{proof}
Branching on a variable with $b$ candidates creates at most $b$ child nodes. Choosing the minimum $b$ minimizes this per-node cost. Under the fail-first principle, this greedily minimizes the expected search tree size.
\end{proof}

\section{Algorithm}

\begin{algorithm}[H]
\caption{Sudoku solver: constraint propagation + backtracking}
\begin{algorithmic}[1]
\REQUIRE Grid with givens
\STATE Initialize $D_{r,c}$ and propagate all givens
\RETURN \textsc{Backtrack}$(D)$
\end{algorithmic}
\end{algorithm}

\begin{algorithm}[H]
\caption{\textsc{Backtrack}$(D)$}
\begin{algorithmic}[1]
\IF{all cells assigned} \RETURN solution \ENDIF
\IF{any $D_{r,c} = \emptyset$} \RETURN FAILURE \ENDIF
\STATE $(r^*,c^*) \gets \arg\min_{(r,c):\;|D_{r,c}|>1} |D_{r,c}|$ \COMMENT{MRV}
\FOR{each $d \in D_{r^*,c^*}$}
    \STATE Save state; assign $d$ to $(r^*,c^*)$; propagate
    \IF{\textsc{Backtrack}$(D) \ne$ FAILURE} \RETURN solution \ENDIF
    \STATE Restore state
\ENDFOR
\RETURN FAILURE
\end{algorithmic}
\end{algorithm}

\section{Complexity Analysis}

\textbf{Time:} Worst case $O(9^{81})$. In practice, propagation reduces well-posed puzzles to trivial search; each of 50 puzzles solves in milliseconds.

\textbf{Space:} $O(81 \times 9) = O(1)$ for domains, plus $O(81)$ stack depth.

\section{Answer}

\[
\boxed{24702}
\]

\end{document}
